% Copyright 2026 Myeongjin Kang. Original prose: CC BY 4.0.
\documentclass[11pt]{article}
\usepackage[T1]{fontenc}
\usepackage{lmodern,microtype,amsmath,amssymb,amsthm}
\usepackage[margin=1.05in]{geometry}
\usepackage[colorlinks=true,linkcolor=blue!50!black,citecolor=blue!50!black,urlcolor=blue!50!black]{hyperref}
\usepackage{xcolor}
\hypersetup{pdftitle={An integral regular summand in every staircase Specht tensor cube},pdfauthor={Myeongjin Kang}}
\newcommand{\Z}{\mathbb Z}
\newcommand{\C}{\mathbb C}
\newcommand{\Sym}{\operatorname{Sym}}
\newcommand{\sgn}{\operatorname{sgn}}
\newcommand{\id}{\operatorname{id}}
\newcommand{\coeff}[2]{[#1]#2}
\newtheorem{theorem}{Theorem}
\newtheorem{lemma}[theorem]{Lemma}
\newtheorem{proposition}[theorem]{Proposition}
\numberwithin{equation}{section}
\title{An integral regular summand in every\texorpdfstring{\\}{ }staircase Specht tensor cube}
\author{Myeongjin Kang}
\date{October 15, 2026}
\begin{document}
\maketitle
\begin{abstract}
For every staircase partition $\rho=(m,m-1,\ldots,1)$ of $N=m(m+1)/2$, we prove that the integral regular module $\Z[S_N]$ is a direct summand of $(S^\rho_{\Z})^{\otimes3}$. Ordered row and column polytabloids give a tensor whose translates form a basis of that summand. A single integral coordinate functional gives an explicit equivariant retraction. The splitting survives base change to every commutative unital ring, including nonflat base change, with no restriction on the characteristic. The proof uses the separating-product mechanism of Bryant and Kov\'acs, Ikenmeyer's triangular palette uniqueness, and a square-deficit calculation from the pinned cyclic-polytabloid construction.
\end{abstract}

\section{Statement and context}\label{sec:intro}
All group actions are left actions, and all tensor products carry the diagonal action. Fix an integer $m\ge1$ and put
\[
 N=\frac{m(m+1)}2,\qquad \rho=(m,m-1,\ldots,1),\qquad
 B=\{(i,j)\in\Z_{\ge1}^2:i+j\le m+1\},\quad G=\Sym(B).
\]
Write $S^\rho_{\Z}$ for the integral Specht lattice in the row-tabloid module, and define $S^\rho_A=A\otimes_{\Z}S^\rho_{\Z}$ for a commutative unital ring $A$. The row and column vectors $v_R,v_C$ below are specific polytabloids in a common word realization of this lattice.

\begin{theorem}\label{thm:main}
For every $m\ge1$, let $L_m=(S^\rho_{\Z})^{\otimes3}$ and $z_m=v_R\otimes v_R\otimes v_C$, with the block orders of Section~\ref{sec:words}. The $\Z[G]$-linear map
\[
 I_m:\Z[G]\longrightarrow L_m,\qquad \delta_g\longmapsto gz_m,
\]
has the explicit $\Z[G]$-linear left inverse $P_m$ in \eqref{eq:P}. Thus the $N!$ translates of $z_m$ form a $\Z$-basis of a split regular summand. For every commutative unital ring $A$, the base-changed maps give a split $A[G]$-linear embedding
\[
 A[G]\longrightarrow (S^\rho_A)^{\otimes_A3},\qquad \delta_g\longmapsto gz_{m,A}.
\]
No flatness or characteristic assumption on $A$ is required.
\end{theorem}

The conclusion concerns the integral and all-characteristic regular summand. The generator and functional are devices for proving that conclusion. Over a splitting field $k$, the regular module decomposes as
\[
 k[G]\cong\bigoplus_{D}P(D)^{\oplus\dim_kD},
\]
where $D$ runs over the simple modules and $P(D)$ is its projective cover. The theorem therefore supplies these multiplicities inside the cube in every characteristic. In characteristic zero this says that each irreducible occurs at least as many times as its dimension.

Harman and Ryba proved complex constituent coverage for the staircase cube \cite[Theorem~1.1]{HR}. Their modular argument already gives every projective indecomposable as a summand in characteristic~2 \cite[Section~3, Remark~3.3]{HR}; that occurrence is not a new consequence here. Regular multiplicities are stronger. The conclusion also applies at odd primes $p\le2m-1$, where the staircase Specht module is not projective: its hook lengths are $2(m+1-i-j)+1$, and the occurrence of a hook of length $p$ makes its dimension fail the divisibility by the order of a Sylow $p$-subgroup required of a projective module. The construction below does not use projectivity of the Specht factor.

The proof combines inherited mechanisms. Bryant--Kov\'acs \cite[Lemma~(ii)--(iii)]{BK} give the regular-summand criterion over a field and the product of separating witnesses. Their proof of the field criterion uses injectivity of the regular group algebra. We instead write and verify an integral retraction directly. Ikenmeyer \cite[Sections~4--5]{Ik} supplies the triangular labeling and its row/column uniqueness. The ordered block construction and the square-deficit step appear in the pinned OpenAI manuscript \cite[Sections~2--3]{OAI}. All needed integral arguments are included here.

\section{Ordered words and polytabloids}\label{sec:words}
Order $B$ lexicographically and let $E=\Z^m$ have ordered basis $e_1,\ldots,e_m$. The basis of $T=E^{\otimes B}$ consists of words $e_a$, indexed by maps $a:B\to\{1,\ldots,m\}$. The position action is
\[
 g e_a=e_{a\circ g^{-1}}.
\]
Tensor regroupings are ordinary and introduce no signs. For an ordered block $D=(d_1,\ldots,d_r)$ set
\[
 \operatorname{alt}_D=\sum_{\sigma\in S_r}\sgn(\sigma)
       \bigotimes_{k=1}^r e_{\sigma(k)},
\]
placing its $k$th factor at $d_k$. Use the exact block orders
\[
 R_i=((i,1),\ldots,(i,m+1-i)),\qquad
 C_j=((1,j),\ldots,(m+1-j,j)),
\]
and define $v_R=\bigotimes_{i=1}^m\operatorname{alt}_{R_i}$ and
$v_C=\bigotimes_{j=1}^m\operatorname{alt}_{C_j}$ in $T$. A word coefficient of either vector is $0$ or $\pm1$; it is nonzero precisely when each block uses its initial-letter palette $\{1,\ldots,r\}$ once.

Here is the integral polytabloid identification, with its signs fixed. A row tabloid of shape $\rho$ with entries in $B$ determines the word assigning letter $k$ to all entries in its $k$th row. This is an equivariant bijection from the tabloid basis to the words of content $\rho$, hence an integral isomorphism onto that coordinate sublattice of $T$. In the tableau $t_R$ put entry $(i,j)$ in cell $(j,i)$; in $t_C$ put it in cell $(i,j)$. Both have shape $\rho$. The ordered columns of $t_R$ are exactly the $R_i$, and those of $t_C$ are the $C_j$. Their initial row words are $j$ and $i$, respectively. The polytabloid
\[
 e_t=\sum_{h\in C_t}\sgn(h)h\{t\}
\]
therefore maps to $v_R$ or $v_C$. Indeed, position permutation replaces the column-letter permutation by its inverse, with the same sign. There is no normalization factor. As every tableau is a relabeling of either fixed tableau, both integral orbit spans equal the same Specht lattice $S^\rho_{\Z}$ inside $T$.

This lattice is finite free, being a subgroup of a finite free abelian group. Tensoring $S^\rho_{\Z}\hookrightarrow T$ successively with free abelian groups embeds $L_m$ into $T^{\otimes3}$. This justifies using ambient word coordinates before base change; it requires no saturation assertion. After base change we use the split maps themselves.

\section{The retraction and its coordinate functional}\label{sec:map}
Let $U$ be free over $\Z$ on letters $f_{r,a}$ with $1\le r\le2m-1$ and $1\le a\le m$. Define
\begin{equation}\label{eq:q}
 q(e_a\otimes e_b\otimes e_c)=
 \begin{cases}f_{a+b-1,a},&a=c,\\0,&a\ne c.\end{cases}
\end{equation}
Apply $q$ at every position after regrouping $T^{\otimes3}$ by positions, obtaining $q^B:T^{\otimes3}\to U^{\otimes B}$. Put
\[
 \mathcal D=\{(r,a):1\le a\le r\le m\}.
\]
Let $W_{\mathcal D}$ be the coordinate sublattice spanned by words using every letter indexed by $\mathcal D$ exactly once and no others. Let $\pi_{\mathcal D}$ be the projection onto it. Both $q^B$ and $\pi_{\mathcal D}$ are $G$-equivariant: the local rule is constant across positions, and content is invariant under position permutation.

Define
\[
 \beta(i,j)=m+2-i-j,\qquad b_0(i,j)=j,\qquad
 w(i,j)=(m+1-i,\beta(i,j)).
\]
The map $w:B\to\mathcal D$ is a bijection, since the row of length $r$ lists $(r,r),(r,r-1),\ldots,(r,1)$. Thus its position orbit is a free transitive basis of $W_{\mathcal D}$, and
\[
 J_m:\Z[G]\xrightarrow{\ \sim\ }W_{\mathcal D},\qquad \delta_g\longmapsto ge_w
\]
is an isomorphism of left modules. Set
\begin{equation}\label{eq:P}
 P_m=J_m^{-1}\pi_{\mathcal D}q^B\big|_{L_m}.
\end{equation}

This construction is exactly a coordinate-functional retraction. Let $\theta^*$ be the coordinate functional of
$\theta=e_\beta\otimes e_{b_0}\otimes e_\beta$ on $T^{\otimes3}$. Then, for every $x\in L_m$,
\begin{equation}\label{eq:coordinate}
 P_m(x)=\sum_{g\in G}\theta^*(g^{-1}x)\,\delta_g.
\end{equation}
To see the equivalence on all basis triples, an output letter $f_{r,a}$ of $q$ has the unique preimage triple $(a,r+1-a,a)$ whenever these are valid input letters. Hence the output word $e_w$ has the unique basis-triple preimage $\theta$. Equivariance gives the unique preimage $g\theta$ of $ge_w$. Its coordinate in $x$ is $\theta^*(g^{-1}x)$, proving \eqref{eq:coordinate}.

Writing $c_R(a)=\coeff{e_a}{v_R}$ and $c_C(a)=\coeff{e_a}{v_C}$, the unit separation identity is consequently
\begin{equation}\label{eq:product}
 \theta^*(gz_m)=c_R(\beta\circ g)c_R(b_0\circ g)c_C(\beta\circ g)
       =\delta_{g,1}.
\end{equation}
These factors realize the separating-product witness of \cite[Lemma~(iii)]{BK}. The alphabet and content projection give a convenient expression for that classical product mechanism. The next section proves the full projected vector identity, which implies \eqref{eq:product} with its exact unit coefficient.

\section{The unique surviving term}\label{sec:proof}
\begin{lemma}[Triangular palette uniqueness, after Ikenmeyer]\label{lem:palette}
A word $a:B\to\{1,\ldots,m\}$ whose letters on row $i$ are exactly $1,\ldots,m+1-i$ and whose letters on column $j$ are exactly $1,\ldots,m+1-j$ must equal $\beta$.
\end{lemma}
\begin{proof}
In the top row, letter $m$ can occur only in column~1, as every later column has smaller maximum. Once its first $k-1$ positions have been filled by $m,m-1,\ldots,m-k+2$, letter $m-k+1$ must occupy column $k$: earlier positions are occupied and later columns have smaller maxima. Thus the top row is $m,m-1,\ldots,1$. Removing it removes the largest allowed letter from each column. Decrease the remaining row indices by one; the palettes are those of the triangle with parameter $m-1$. Induction, starting at $m=1$, proves the claim. Conversely $\beta$ has all the required palettes. This is the row/column uniqueness used in \cite[Section~5, equation~(5.2)]{Ik}.
\end{proof}

\begin{proposition}\label{prop:unit}
With the stated block orders, $\pi_{\mathcal D}q^B(z_m)=e_w$ over $\Z$.
\end{proposition}
\begin{proof}
An expansion term of $z_m$ is a basis triple $(a,b,c)$. The words $a,b$ independently have the initial-letter palette on each row, and $c$ has that palette on each column. Each triple appears once, with the product of its three block-permutation signs.

Consider any term retained by both maps. The local rule forces $a=c$. Write $d=a+b-1$ pointwise. The output contains each $(r,s)\in\mathcal D$ once, so
\begin{equation}\label{eq:global-square}
 \sum_{x\in B}d(x)^2=\sum_{r=1}^m r^3.
\end{equation}
On a row of length $r$, both $a$ and $b$ are permutations of $1,\ldots,r$. Therefore
\begin{equation}\label{eq:deficit}
 \sum_{\text{row}}d=r^2,\qquad
 \sum_{\text{row}}(d-r)^2=\sum_{\text{row}}d^2-r^3\ge0.
\end{equation}
Summing over the rows and using \eqref{eq:global-square} makes the sum of these nonnegative integer squares zero. Consequently $d=r$ at every position in the row of length $r$, and
\[
 b(i,j)=m+2-i-a(i,j).
\]
This is the square-deficit calculation of \cite[Section~3, \texttt{dom:square-minimum}]{OAI}. It concerns the integer labels of individual support terms, not an ordering or positivity property in a coefficient ring.

The common word $a=c$ has both the row and column palettes. Lemma~\ref{lem:palette} gives $a=c=\beta$, and then $b=b_0$. Conversely this triple satisfies every block palette and maps to $e_w$. It is therefore the only surviving expansion term. With
\[
 H=\sum_{r=1}^m\binom r2,
\]
its first layer reverses each row, its second layer is increasing on each row, and its third layer reverses each column. The coefficient is exactly
\[
 (-1)^H\cdot1\cdot(-1)^H=+1.
\]
This proves the full vector identity, including $m=1$.
\end{proof}

\begin{proof}[Proof of Theorem~\ref{thm:main}]
Equivariance and Proposition~\ref{prop:unit} give $P_m(gz_m)=\delta_g$ for every $g\in G$, so $P_mI_m=\id$. Explicitly,
\[
 L_m=I_m(\Z[G])\oplus\ker P_m,\qquad
 x=I_mP_mx+(x-I_mP_mx).
\]
For any commutative unital ring $A$, functoriality of tensor product gives
\[
 (1_A\otimes P_m)(1_A\otimes I_m)=\id_{A\otimes_{\Z}\Z[G]}.
\]
Under the canonical identifications $A\otimes_{\Z}\Z[G]=A[G]$ and
$A\otimes_{\Z}L_m=(S^\rho_A)^{\otimes_A3}$, the embedded basis vectors are the stated translates of $z_{m,A}$. Tensoring preserves this equality of split maps without requiring $A$ to be flat. In particular no factorial is inverted and no field-level injectivity theorem is used over $\Z$.
\end{proof}

\section{Relation to earlier conclusions and verification}\label{sec:boundary}
The complex regular conclusion already follows from the earlier project note \cite[Section~2.1]{I21} and Harman--Ryba \cite[Theorem~2.4]{HR}. The former identifies $\C[G](v_R\otimes v_R)$ with $M^\rho$, using the pair-letter projection of \cite{OAI}, whose scalar is
\[
 (-1)^H\prod_{r=1}^m r!.
\]
The latter puts $M^{C(\rho)}=M^{(1^N)}=\C[G]$ inside $M^\rho\otimes S^\rho$. Semisimplicity gives complex regular containment in the cube. That comparison does not give an integral splitting, since the scalar need not be a unit. The present third factor and refined letters give the unit identity of Proposition~\ref{prop:unit}.

For context, we use OpenAI revision \texttt{fd4aeeb2ee4f} (the full pin is recorded in the references). The pinned source stipulates the stronger complex Saxl square theorem \cite[Theorem~3.1, \texttt{thm:cyclic-saxl}]{OAI}; its companion \cite{Universal} treats universal complex squares. Those results are accepted as stipulated source context here. The integral proof above does not use their complex coverage conclusions. The characteristic-2 projective occurrence in \cite{HR} follows the modular approach of Bessenrodt--Bowman--Sutton \cite[Section~5]{BBS}; their strengthened conjecture concerns the square. The square-splitting identity in \cite[Theorem~3.7]{BB} concerns a different multiplicity question. We claim neither those results nor the classical separating criterion, triangular uniqueness, or square-deficit mechanism as new. The advance asserted in this note is the integral and every-ring regular-summand package. No worldwide priority claim is made.

The accompanying code copies two previously independent finite controls with their provenance preserved. One constructs polytabloids and checks all $1+6+720=727$ orbit elements for $m\le3$, including equivariance and sign mutations. The other tests the product-coordinate condition for $m\le5$ by exhaustive backtracking with literal block-support pruning, and includes an unpruned $m=3$ check and a characteristic-2 obstruction to splitting the intermediate permutation module from the square at $m=2$. These calculations are reproducibility support, not premises of the all-$m$ proof. No full Lean formalization is supplied or claimed. Exact source interfaces, hashes, build instructions, and the formal-coverage status accompany the source.

\paragraph{Rights.}
\copyright\ 2026 Myeongjin Kang. Original prose is licensed under \href{https://creativecommons.org/licenses/by/4.0/}{CC BY 4.0}; accompanying original code is MIT-licensed. Third-party works retain their own rights. Primary-paper reading copies are excluded from the redistributable package.

\input{references.tex}
\end{document}
